\documentclass{template}

% --- PACKAGES ---
% \usepackage[margin=1in, headheight=14pt, top=1.2in]{geometry}
% \usepackage{amsmath,amssymb}
% \usepackage{graphicx}
% \usepackage{enumitem}
% \usepackage{hyperref}
% \usepackage{fancyhdr}
% \usepackage{float}
\usepackage{xcolor}
\usepackage{mdframed}
\usepackage{etoolbox}
\usepackage{pdfpages}

\graphicspath{{attachments/}}

% --- TOGGLE FOR ANSWERS ---
% Set to \toggletrue{showanswers} to display answers.
% Set to \togglefalse{showanswers} to hide answers.
\newtoggle{showanswers}
\toggletrue{showanswers}

% --- CUSTOM SOLUTION ENVIRONMENT ---
\newmdenv[
  backgroundcolor=white,
  linecolor=blue!60!black,
  linewidth=1.5pt,
  topline=false,
  rightline=false,
  bottomline=false,
  leftmargin=2pt,
  innerleftmargin=10pt,
  innertopmargin=6pt,
  innerbottommargin=6pt,
  skipabove=8pt,
  skipbelow=8pt
]{solutionbox}

\newcommand{\solution}[1]{%
\iftoggle{showanswers}{%
\begin{solutionbox}
        #1
\end{solutionbox}%
  }{}%
}

% \usepackage[]{opensans} % Loads Open Sans font family
% \newcommand{\solution}[1]{%
%   \iftoggle{showanswers}{%
%     \begin{solutionbox}
%         \fontsize{10}{12}\opensans #1
%     \end{solutionbox}%
%   }{}%
% }

\addbibresource{/Users/nathalieredick/repos/washu/Library.bib}


% --- HEADER & FOOTER CONFIGURATION ---
\pagestyle{fancy}
\fancyhf{}
\lhead{\textbf{EEPS 5534: Interior of the Earth}}
\rhead{\textbf{Problem Set 2}}
\cfoot{\thepage}


\title{\textbf{Problem Set 2}}
\author{Nathalie Redick}
\date{Due: Oct. 09, 2026}

\begin{document}

\maketitle
\thispagestyle{fancy}

\begin{enumerate}[label=\textbf{\arabic*.}]

% ------------------------------------------------------------
% QUESTION 1
% ------------------------------------------------------------
\item Shear wave splitting observations of seismic anisotropy from vertically propagating SKS waves commonly report a splitting time (time difference between the fast and slow S wave) on the order of 1 second. Discuss the interpretation of such observations, making reasonable estimates for the seismic velocities and anisotropy \dots

\begin{enumerate}[label=\Alph*)]
\item Crustal rocks are assumed to show anisotropy on the order of 1\%. Make a simple calculation showing why the observed splitting cannot be caused by crustal rocks.
\solution{
% wip ADD YOUR ANSWER HERE
    Given that the average composition of continental crust is granitic, we can assume a shear wave velocity of the crust of 3.67 km/s (given by question 2). Using the following equation for delay time, as defined by \textcite{silverSeismicAnisotropyContinents1996,silverShearWaveSplitting1991}:

    \begin{equation}
        \delta t = L \frac{\delta\hat{\beta}}{\beta_0}
    \end{equation}

    Where:

    \begin{itemize}
        \item $L$ = path length through the crust
        \item $\delta\hat{\beta}$ = dimensionless anisotropy
        \item $\beta_0$ = isotropically averaged shear velocity
    \end{itemize}

    The expected delay time for continental crust, assuming an average thickness of 40~km \parencite{mooneyCrustLithosphericStructure2015}, can be estimated as:

    \begin{align*}
        \delta t &= \qty{40}{\km} \frac{0.01}{\qty{3.67}{\km\per\s}}\\
        &= \qty{0.1}{\s} \ll\, \qty{1}{\s}
    \end{align*}

    Even for exceptionally thick regions of the crust estimated at \textasciitilde70~km \parencite{mooneyCrustLithosphericStructure2015}, the expected delay time is only:

    \begin{align*}
        \delta t &= \qty{70}{\km} \frac{0.01}{\qty{3.67}{\km\per\s}}\\
        &= \qty{0.2}{\s} \ll\, \qty{1}{\s}
    \end{align*}

    This is consistent with delay time estimates of 0.1--0.2~s by \textcite{savageSeismicAnisotropyMantle1999,silverSeismicAnisotropyContinents1996}. Even for the thickest regions of the crust, the expected anisotropy only accounts for a small fraction of the observed anisotropy, therefore, the observed splitting cannot be caused by only crustal rocks. 
}

\item Now calculate and make a graph showing the trade-off between the inferred thickness of the upper mantle anisotropic layer and the estimated percent anisotropy of the mantle rocks, assuming a single anisotropic layer. How might one be able to constrain this trade off?
\solution{
% todo ADD YOUR ANSWER HERE

}
\end{enumerate}

% ------------------------------------------------------------
% QUESTION 2
% ------------------------------------------------------------
\item Samples of granite and peridotite are examined in the laboratory. The P-wave velocities of the rocks are 6.25 km/s and 8.3 km/s, the S-wave velocities are 3.67 km/s and 4.73 km/s, and the densities are 2650 kg/m$^3$ and 3300 kg/m$^3$.

\begin{enumerate}[label=\Alph*)]
\item Find the Lam\'e parameters ($\lambda$ and $\mu$) and bulk modulus. Be sure to give the correct units. Are these rocks well approximated as Poisson solids?
\solution{
% todo ADD YOUR ANSWER HERE

}

\item The bulk modulus is the ratio of the pressure to the volume change of a material. Estimate the change in the volumes and densities of the two rocks if they are buried to a depth of 40 km.
\solution{
% todo ADD YOUR ANSWER HERE

}
\end{enumerate}

% ------------------------------------------------------------
% QUESTION 3
% ------------------------------------------------------------
\item Near the earth's surface, the soil is generally in equilibrium with the mean annual temperature (except within a couple meters of the surface, where there are seasonal and diurnal variations). Recent evidence suggests that some climatic changes are extremely rapid. Assume that during the last deglaciation in Scotland, the mean annual temperature in Scotland suddenly changed from 0 to 10~$^\circ$C (for the sake of this problem you can approximate the change as instantaneous, and ignore geothermal heat flux). Calculate temperature as a function of time using the formula given in class and plot for the last 15,000 years at depths of 50~m, 200~m and 1~km. Assume that the conductivity of the soil is 1.5~W/m and the heat capacity is 700~J/(kg-K). Comment on the usefulness of soil temperature profiles for reconstructing the climate record.
\solution{
% todo ADD YOUR ANSWER HERE

}

% ------------------------------------------------------------
% QUESTION 4
% ------------------------------------------------------------
\item Continental cratons have deep, high seismic velocity lithospheric ``roots'' with velocity anomalies of 5--8\% relative to the surrounding mantle. Now test the idea that these ``roots'' are due to the cooling of the mantle lithosphere. The roots are observed down to depths of about 250~km for old cratons (age $\sim$2.5~Ga).

\begin{enumerate}[label=\Alph*)]
\item Calculate and plot the temperature at 250~km depth for the halfspace cooling model as a function of time.
\solution{
% todo ADD YOUR ANSWER HERE

}

\item Using the temperature derivative of seismic velocity and the temperature you calculated in part A), calculate the seismic velocity anomaly expected at 250~km depth for a cooling time of 2.5~Ga. Assume that the temperature of the surrounding mantle does not change. Can the ``roots'' be due to cooling?
\solution{
% todo ADD YOUR ANSWER HERE

}
\end{enumerate}

% ------------------------------------------------------------
% QUESTION 5
% ------------------------------------------------------------
\item \textbf{Fitting concentration profiles with experimentally determined diffusion parameters, and understanding volcanic timescales through simple diffusion chronometry.}

Timescales of magmatic processes relevant to igneous rocks can be obtained by modeling the diffusive relaxation of chemical profiles in minerals. For this problem, we will use data from specially designed experiments in order to determine the general Arrhenius parameters for Fe-Mg inter-diffusion in olivine. Then these parameters will be used to fit real data from an olivine crystal from Mt.\ Etna. The goal is to determine the high-T storage time of the olivine grain in order to understand how much time passed between crystallization and eruption.

\begin{enumerate}[label=\textit{Part \arabic*:}, leftmargin=*]

\item Using the data provided from Chakraborty (1997) \textit{J.\ Geophys.\ Res.} The file \\\texttt{problem\_set\_2\_2026.xls} contains concentration profiles from 4 experiments that used gem-quality single crystals of olivine of different compositions in a diffusion couple experiment to generate 1-D concentration gradients. The 4 experiments are carried out at different temperatures and for different durations. With this information, you can solve for the diffusion coefficient for interdiffusion of Fe and Mg. Solve for the $D$ value for the 4 different experiments, then plot this data on a $D$ vs $1/T$ plot. Do a linear regression fit to the data to determine the $D_0$ and slope which is the necessary information to make a generalized Arrhenius law that you can use to calculate diffusive timescales in olivine. For this problem, realize that the $f_{\mathrm{O}_2}$ is fixed, and you can just assume that the Mt.\ Etna olivine is at the same $f_{\mathrm{O}_2}$, so can be ignored in this problem set. Also do not worry about taking into account anisotropic effects of diffusion. I have provided a tab in excel where you can practice with the error function fitting and plotting. Please feel free to do the fits and regressions in any program such as jupyter, matlab, python, etc.

In order to solve for the $D$ value, you will need make a complimentary error function fit to the data using the following equation which we have gone over in class:
\[
C(x,t) = \tfrac{1}{2}\left(C_L - C_R\right)\,\mathrm{erfc}\!\left(\frac{x}{2\sqrt{Dt}}\right) + C_R
\]
where $C$ refers to concentration in molar units (using mol of Fo content or Fa content varying from 0--1 is fine). $C_L$ is the concentration on the left hand side, $C_R$ the right hand side of the diffusion couple. You will also need to redefine the distance axis so that `0' occurs at the midpoint of the trend. Adjust the $D$ value empirically to achieve the best fit. It would be more accurate to linearize the profile and do a least squares regression, if you want to go the extra mile, but for the sake of this homework, eyeball empirical fits are just fine! Make sure to convert all units to SI units, m, s, K. You can solve for $D_{\mathrm{Fe}}$ or $D_{\mathrm{Mg}}$, but be consistent across all 4 experiments.

Now once you have fit all the experimental diffusion profiles, collect the data together in the `Diffusivities\_all' tab of the spreadsheet. Make an Arrhenius plot, fit the data with a least squares linear regression with the following equation style:
\[
D = D_0 \exp\!\left(-\frac{Q}{RT}\right)
\]
where $D_0$ is the pre-exponential constant (diffusivity at infinite $T$), and $R$ is the gas constant and $Q$ is the activation energy. You can do a simple least squares fit here, but a more accurate way would be to take into account uncertainty on the points (obtained from your fits if you did a linearized fitting to the experiments) and then do a weighted least squares regression. But this is not necessary for this problem set.
\solution{
% todo ADD YOUR ANSWER HERE

}

\item \textit{Mt.\ Etna:} On the `natural olivine' tab you have electron microprobe data of a concentration profile in an olivine from Mt.\ Etna. The dominant feature of this olivine is that it is reversely zoned. You interpret this inverse zoning as indicating a discreet magma mixing event with a more mafic magma, inducing a thick rim of Mg-rich olivine to grow. Initially this is a sharp boundary, but is relaxed by chemical diffusion.

Using the diffusion parameters from \textit{Part 1}, make a best fit for the Etna data. The temperature for this eruption was independently determined to be 1150~$^\circ$C. Determine a timescale of mixing (time from the mixing event when the diffusive clock started to eruption when the diffusion clock effectively stops). You can ignore $f_{\mathrm{O}_2}$, pressure, and crystallographic effects for this calculation.

Note, the thin normally-zoned rim at the very edge of the profile. What process(es) could this be related to?
\solution{
% todo ADD YOUR ANSWER HERE

}

\end{enumerate}

% ------------------------------------------------------------
% QUESTION 6
% ------------------------------------------------------------
\item Geochemists trying to balance Earth's heat budget face a puzzle regarding heat-producing radioactive elements (U, Th, and K). Radiogenic heat is generated primarily by the decay of long-lived radioactive isotopes: Uranium-238 ($^{238}$U), Thorium-232 ($^{232}$Th), and Potassium-40 ($^{40}$K). Because these are highly incompatible elements, they are heavily concentrated in the crust during partial melting.

In this problem, we will make some simple assumptions and try to calculate the heat budget being produced in the stable continental crust of the Earth, and the implications for its overall composition.

\textbf{Dataset and parameters:}

Assume a felsic upper crustal rock with density $\rho = 2700$~kg/m$^3$, $C_U = 2.8~\mu$g/g, $C_{Th} = 10.7~\mu$g/g, and $C_K = 2.5$~wt\%.

Continental crust mean heat flow within stable continental shields: $\sim$60~mW/m$^2$.

\textbf{Theoretical heat production equations:}

The specific radiogenic heat production rate per unit mass of rock is estimated as follows:
\begin{align*}
\text{U} &= 9.52\times10^{-5}~\text{W/kg}\\
\text{Th} &= 2.56\times10^{-5}~\text{W/kg}\\
\text{K} &= 3.48\times10^{-9}~\text{W/kg}
\end{align*}

The per volume heat production ($A$, in W/m$^3$) is calculated using the Rybach equation, multiplying density ($\rho$) by the heat production by mass:
\[
A = \rho \times (\text{total W/kg of radiogenic heat production})
\]
where $\rho$ = rock density (kg/m$^3$). To scale volumetric heat production ($A$) up to surface heat flow ($q_s$, in W/m$^2$), we will use the linear heat flow model from Roy et al., 1968:
\[
q_s = q_r + D \times A
\]
where $q_r$ is the reduced heat flow from deeper mantle/core sources ($\sim 15\times10^{-3}$~W/m$^2$) and $D$ is the characteristic crustal thickness scale (m) over which heat-producing elements are concentrated.

\begin{enumerate}[label=\alph*)]
\item Assuming an upper crustal thickness of 10~km, and using the provided abundances of heat-producing elements in the upper crust, what is the surface heat flow ($q_s$) from the upper crust and mantle/core sources, and how does this compare to the total average value of 53~mW/m$^2$?
\solution{
% todo ADD YOUR ANSWER HERE

}

\item Assuming the average stable continental crust has a thickness of 30~km, the total surface heat flow is 53~mW/m$^2$, and the density of the lower crust is 2900~kg/m$^3$, what is the concentration of U, Th, and K in the lower crust? (Assume the K/U and Th/U are the same for upper and lower continental crust.) How does this compare with the concentration in the upper crust? How does it compare to the average n-MORB (see Table 3 in Gale et al.\ 2013, \textit{G$^3$})?
\solution{
% todo ADD YOUR ANSWER HERE

}

\item If the continental crust holds a massive share of the planet's incompatible elements despite accounting for $<$\,1\% of BSE mass, what does this imply about the depletion state of the upper mantle where MORBs (Mid-Ocean Ridge Basalts) are generated?
\solution{
% todo ADD YOUR ANSWER HERE

}
\end{enumerate}

\end{enumerate}

% \section*{Acknowledgments}
% Some LaTeX typesetting/formatting, conversion of table for question 5 to CSV format, and code clean-up assistance provided by Google Gemini (September 2026 version).

\singlespacing
\printbibliography[title={References}]

% \newpage

\end{document}